% ===================== Chapter 14: NumPy (ARC500) =====================
% Revised for architecture students without vector/matrix background
% Requires your project preamble to define:
% \h{...}, \hh{...}, \hhh{...}
% \code{<label>}{python}{<code>}{<caption>}
% All Python examples below use EXACTLY 4 spaces per indent (no tabs).

\h{NumPy (Numerical Python) for Computing}

Imagine you're working on a 50-story tower with 20 apartments per floor. That's 1,000 units, each with data about area, orientation, window ratios, and more. Using regular Python lists and loops to calculate total areas or average sizes would require writing many repetitive lines of code. NumPy lets you perform these calculations on all 1,000 units at once with a single command---like using Excel formulas but much more powerful.

\emph{NumPy} is a Python library that makes working with large sets of numbers incredibly fast and simple. Think of it as a supercharged spreadsheet that lives in your Python code. While Python lists are like filing cabinets where you can store anything (numbers, text, images), NumPy arrays are like organized grids designed specifically for numbers---making calculations thousands of times faster.

\begin{tcolorbox}[
    title=NumPy vs Traditional Approach,
    colframe=black,
    colback = white,
    arc=2mm
]
\textbf{Task: Calculate total area for 1000 apartments}
\begin{lstlisting}[basicstyle=\ttfamily\footnotesize,breaklines=true,columns=fullflexible,keepspaces=true,xleftmargin=0pt]
Traditional Python (with loops):         NumPy (vectorized):
                                         
areas = []                               areas = np.array([...1000 values...])
for i in range(1000):    
    area = width[i] * length[i]  <--+   total = areas.sum()  
    areas.append(area)              |    
                                    |    Done in 1 line!
total = 0                           |    0.001 seconds
for area in areas:                  |    
    total += area                <--+    
                                         
10+ lines of code                       
0.1 seconds                             

Speed comparison for different building sizes:
Units:     10    100    1,000   10,000   100,000
Lists:   0.001s  0.01s   0.1s     1s       10s
NumPy:   0.001s  0.001s  0.001s   0.002s   0.02s
                 +-------+--------+
              NumPy stays fast even with huge datasets!

\end{lstlisting}
\end{tcolorbox}

\hh{Applications}

Before diving into code, let's see where you'll use NumPy in practice:

\begin{itemize}
\item \textbf{Space planning:} Calculate areas for hundreds of rooms, find which meet code requirements, sum by floor or building wing
\item \textbf{Facade design:} Generate coordinates for thousands of panels, apply sun angles to each, calculate shading percentages
\item \textbf{Structural grids:} Define column locations, beam spans, and load distributions across entire buildings
\item \textbf{Environmental analysis:} Process hourly temperature data for a year (8,760 values), find averages, peaks, and patterns
\item \textbf{Cost estimation:} Apply unit costs to quantities, sum by trade or building zone, apply regional factors
\item \textbf{Integration with design tools:} Exchange data with Rhino/Grasshopper, export to Excel for consultants, prepare data for visualization
\end{itemize}

\begin{tcolorbox}[
    title=Focus of this chapter:,
    colframe=orange,
    colback = white,
    arc=2mm
]
Think of this chapter as learning to use a powerful calculator that works on entire building datasets at once:
\begin{itemize}
\item Create number grids (arrays) representing rooms, floors, or facades
\item Read array properties: size, dimensions, data types (like checking spreadsheet dimensions)
\item Select specific floors, rooms, or zones (like filtering in Excel)
\item Apply formulas to entire datasets at once (like Excel's fill-down, but more powerful)
\item Calculate totals, averages, and other statistics by floor, wing, or entire building
\item Reshape data from lists to grids and back (like pivot tables)
\item Use built-in math functions for common architectural calculations
\end{itemize}
\end{tcolorbox}

% ====================================================================
% PART I -- NUMPY BASICS
% ====================================================================

\hh{Getting started (install and import)}

NumPy isn't included with Python by default---you need to install it once on your computer, then import it in any script where you want to use it. The architecture community (and everyone else) uses the nickname 'np' for NumPy. Think of it like how we say "CAD" instead of "Computer-Aided Design."

\code{c14_setup}{python}{
    # STEP 1: Install NumPy (do this ONCE in your terminal/command prompt)
    # Mac: open Terminal
    # Windows: open Command Prompt or PowerShell
    # Type: pip install numpy
    # Press Enter and wait for installation to complete

    # STEP 2: In your Python script, import NumPy
    import numpy as np  # 'np' is the universal nickname
    
    # Check it worked
    print("NumPy version:", np.__version__)
    # Outcome:
    # NumPy version: 2.0.x (or similar)
}{Setting up NumPy on your computer}

\hh{Arrays vs lists: Understanding the difference}

Let's start with a simple architectural example to see why NumPy arrays are special:

\code{c14_why_numpy}{python}{
    import numpy as np
    import time  # To measure performance

    # SCENARIO: 5 studio apartments with widths and lengths in meters
    
    # Method 1: Python lists (the slow, tedious way)
    widths_list = [3.5, 4.0, 5.0, 3.8, 4.2]   # regular Python list
    lengths_list = [8.0, 7.5, 6.0, 8.0, 7.0]  # regular Python list
    
    # To calculate areas with lists, we need a loop:
    areas_list = []
    for i in range(5):
        area = widths_list[i] * lengths_list[i]
        areas_list.append(area)
    print("Areas (using lists):", areas_list)
    
    # Method 2: NumPy arrays (the fast, elegant way)
    widths = np.array([3.5, 4.0, 5.0, 3.8, 4.2])   # NumPy array
    lengths = np.array([8.0, 7.5, 6.0, 8.0, 7.0])  # NumPy array
    
    # With NumPy, multiply entire arrays at once!
    areas = widths * lengths  # That's it! No loop needed
    print("Areas (using NumPy):", areas)
    
    # Performance comparison with 10,000 rooms
    print("\n=== Performance Comparison ===")
    large_widths_list = list(np.random.uniform(3, 6, 10000))
    large_lengths_list = list(np.random.uniform(6, 10, 10000))
    
    # Time the list method
    start = time.time()
    areas_list = []
    for i in range(10000):
        areas_list.append(large_widths_list[i] * large_lengths_list[i])
    list_time = time.time() - start
    
    # Time the NumPy method
    large_widths = np.array(large_widths_list)
    large_lengths = np.array(large_lengths_list)
    start = time.time()
    areas_numpy = large_widths * large_lengths
    numpy_time = time.time() - start
    
    print(f"List method: {list_time:.3f} seconds")
    print(f"NumPy method: {numpy_time:.4f} seconds")
    print(f"NumPy is {list_time/numpy_time:.0f}x faster!")
    
    # Outcome:
    # Areas (using lists): [28.0, 30.0, 30.0, 30.4, 29.4]
    # Areas (using NumPy): [28.  30.  30.  30.4 29.4]
    # 
    # === Performance Comparison ===
    # List method: 0.003 seconds
    # NumPy method: 0.0001 seconds
    # NumPy is 30x faster!
}{NumPy eliminates loops and runs much faster}

\hh{Visual understanding: how arrays represent buildings}

\begin{tcolorbox}[
    title=Visualizing arrays as buildings,
    colframe=purple,
    colback = white,
    arc=2mm,
    breakable
]
\textbf{1D array = single floor (linear arrangement)}

Understanding how NumPy stores data differently from Python lists is crucial. The diagram below shows the memory layout difference, which explains why NumPy is so much faster.

\begin{lstlisting}[basicstyle=\ttfamily\footnotesize,breaklines=true,columns=fullflexible,keepspaces=true,xleftmargin=0pt]
Python List vs NumPy Array:

Python List (slow):     NumPy Array (fast):
    +----+                 +----+----+----+----+
    | 45 | --+             | 45 | 52 | 48 | 55 |
    +----+   |             +----+----+----+----+
    +----+   |              Unit: 0   1   2   3
    | 52 | --+ scattered    
    +----+   | in memory    All values stored
    +----+   |              continuously in memory
    | 48 | --+              = much faster access!
    +----+   |
    +----+   |
    | 55 | --+
    +----+

apartments = np.array([45, 52, 48, 55])  # Areas in m^2
\end{lstlisting}

\textbf{What this diagram teaches:}
\begin{itemize}
\item \textbf{Python list (left):} Each value is stored in a different location in memory. The arrows show that Python must "jump around" memory to find each value, which is slow.
\item \textbf{NumPy array (right):} All values are stored in one continuous block. The computer can read them in one sweep, like reading a sentence rather than hunting for each letter.
\item \textbf{The unit numbers (0, 1, 2, 3):} Show how we access elements. In both cases, we start counting at 0, not 1. This is universal in programming.
\item \textbf{Why it matters:} When processing thousands of apartments, these memory jumps in Python lists become a major bottleneck.
\end{itemize}
\vspace{2em}

\textbf{2D array = multi-story building}

Now let's see how NumPy represents a multi-floor building. This is like going from a single row in Excel to a full spreadsheet.

\begin{lstlisting}[basicstyle=\ttfamily\footnotesize,breaklines=true,columns=fullflexible,keepspaces=true,xleftmargin=0pt]
Building as 2D Array:

                  Unit 0   Unit 1   Unit 2   Unit 3
                 +--------+--------+--------+--------+
    Floor 2 -->  |   44   |   51   |   49   |   56   |
                 +--------+--------+--------+--------+
    Floor 1 -->  |   46   |   53   |   47   |   54   |
                 +--------+--------+--------+--------+
    Floor 0 -->  |   45   |   52   |   48   |   55   |
    (Ground)     +--------+--------+--------+--------+
                     |        |        |        |
                  axis=0   axis=0   axis=0   axis=0
                  (down)   (down)   (down)   (down)
    
    <---------- axis=1 (across) ---------->

building = np.array([[45, 52, 48, 55],  # Floor 0
                     [46, 53, 47, 54],  # Floor 1
                     [44, 51, 49, 56]]) # Floor 2

Shape: (3, 4) = 3 floors x 4 units
\end{lstlisting}

\textbf{Understanding this 2D representation:}
\begin{itemize}
\item \textbf{The grid:} Each cell contains an apartment's area in square meters. This is exactly like a spreadsheet where rows are floors and columns are unit positions.
\item \textbf{Floor numbering:} We start at Floor 0 (ground floor). This matches how Python counts and how many countries number floors.
\item \textbf{The axes labels:} These are critical! Axis=0 goes vertically (through floors), axis=1 goes horizontally (across units). You'll use these constantly.
\item \textbf{Shape (3, 4):} This notation means 3 rows by 4 columns. Always read shapes as (rows, columns).
\item \textbf{The code below:} Shows how to create this 2D array. Notice the double square brackets---the outer brackets contain the array, inner brackets contain each floor.
\end{itemize}
\vspace{2em}

\textbf{Understanding axes}

This is perhaps the most important concept in NumPy. Axes determine the direction of operations like sums and averages.

\begin{lstlisting}[basicstyle=\ttfamily\footnotesize,breaklines=true,columns=fullflexible,keepspaces=true,xleftmargin=0pt]
Key Concept: Axes determine direction of operations

        axis=1 (horizontal) -->
        +----+----+----+----+
axis=0  | 45 | 52 | 48 | 55 | } sum(axis=1) = 200
  |     +----+----+----+----+     (sum across row)
(vert)  | 46 | 53 | 47 | 54 | } sum(axis=1) = 200
  |     +----+----+----+----+
  v     | 44 | 51 | 49 | 56 | } sum(axis=1) = 200
        +----+----+----+----+
          |    |    |    |
        135  156  144  165 = sum(axis=0)
                              (sum down column)
\end{lstlisting}

\textbf{Critical axis concepts:}
\begin{itemize}
\item \textbf{axis=0 (vertical):} Operations along this axis go DOWN through rows. When you sum along axis=0, you're adding up all floors for each unit position.
\item \textbf{axis=1 (horizontal):} Operations along this axis go ACROSS columns. When you sum along axis=1, you're adding up all units on each floor.
\item \textbf{Memory trick:} axis=0 is like dropping down an elevator shaft, axis=1 is like walking down a hallway.
\item \textbf{The sums shown:} 200 is the total area per floor (adding horizontally), while 135, 156, etc. are totals for each unit position through all floors (adding vertically).
\item \textbf{Why this matters:} You'll constantly use axes to answer questions like "What's the total area per floor?" (axis=1) or "What's the average size of corner units?" (axis=0).
\end{itemize}
\vspace{2em}

\textbf{Reshaping = reorganizing units}

Reshaping is like taking the same apartments and arranging them differently---from a long corridor to a tower, or vice versa.

\begin{lstlisting}[basicstyle=\ttfamily\footnotesize,breaklines=true,columns=fullflexible,keepspaces=true,xleftmargin=0pt]
From Linear to Grid:

Original (12 units in a line):
+---+---+---+---+---+---+---+---+---+----+----+----+
| 1 | 2 | 3 | 4 | 5 | 6 | 7 | 8 | 9 | 10 | 11 | 12 |
+---+---+---+---+---+---+---+---+---+----+----+----+
Shape: (12,)  <-- Note the comma! This means 1D array

reshape(3, 4) --> 3 floors x 4 units:
+---+---+---+----+
| 1 | 2 | 3 | 4  | Floor 0
+---+---+---+----+
| 5 | 6 | 7 | 8  | Floor 1
+---+---+---+----+
| 9 |10 |11 | 12 | Floor 2
+---+---+---+----+
Shape: (3, 4)

reshape(2, 6) --> 2 wings x 6 units:
+---+---+---+---+----+----+
| 1 | 2 | 3 | 4 | 5  | 6  | Wing A
+---+---+---+---+----+----+
| 7 | 8 | 9 |10 | 11 | 12 | Wing B
+---+---+---+---+----+----+
Shape: (2, 6)
\end{lstlisting}

\textbf{Understanding reshaping:}
\begin{itemize}
\item \textbf{The numbers 1-12:} Represent apartment IDs or any sequential data. Notice how the same 12 values can be arranged in different configurations.
\item \textbf{Shape notation:} (12,) with a comma means a 1D array with 12 elements. (3,4) means 2D with 3 rows and 4 columns.
\item \textbf{Reading order:} NumPy fills the new shape by reading left-to-right, top-to-bottom, just like reading English text.
\item \textbf{The constraint:} Total elements must stay the same. You can't reshape 12 elements into a 3x5 grid (that would need 15 elements).
\item \textbf{Real-world use:} You might receive data as a long list from a sensor or database, then reshape it to match your building's actual floor plan.
\item \textbf{Different interpretations:} The same reshape can represent different things---3x4 could be 3 floors with 4 units each, or 3 building sections with 4 data points each.
\end{itemize}
\end{tcolorbox}

\hh{Creating arrays: Your first building blocks}

Arrays are like organized tables of numbers. Let's explore different ways to create them, using architectural examples:

\code{c14_create_arrays_detailed}{python}{
    import numpy as np
    
    # METHOD 1: Convert existing data into an array
    # Like importing room areas from a spreadsheet
    room_areas_m2 = np.array([45.2, 32.1, 28.5, 51.3, 38.7])  # ALWAYS label units!
    print("Room areas (m^2):", room_areas_m2)
    
    # METHOD 2: Create a sequence with arange (like "array range")
    # Perfect for regular grids or floor numbers
    floors = np.arange(0, 10)        # floors 0 through 9
    print("Floor numbers:", floors)
    
    # Column grid at 6m spacing
    column_positions_m = np.arange(0, 31, 6)  # 0, 6, 12, 18, 24, 30 meters
    print("Column positions (m):", column_positions_m)
    
    # METHOD 3: Create evenly spaced points with linspace
    # Useful for dividing a facade into equal parts
    # Divide a 30m facade into 5 equal sections (6 division points)
    facade_divisions_m = np.linspace(0, 30, 6)
    print("Facade divisions (m):", facade_divisions_m)
    
    # METHOD 4: Create arrays filled with the same value
    # Like initializing a design with standard values
    
    # 5 rooms, all starting with 0 occupancy
    occupancy = np.zeros(5)
    print("Initial occupancy:", occupancy)
    
    # Standard ceiling height for 8 rooms
    ceiling_heights_m = np.full(8, 3.2)  # 3.2m for all rooms
    print("Ceiling heights (m):", ceiling_heights_m)
    
    # METHOD 5: Arrays with construction constraints
    # Snap to 0.1m construction grid
    raw_dimensions = np.array([3.142, 4.678, 5.823, 2.456])
    snapped_dimensions = np.round(raw_dimensions, 1)
    print("\nConstruction grid (0.1m):")
    print("Raw:", raw_dimensions)
    print("Snapped:", snapped_dimensions)
    
    # Outcome:
    # Room areas (m^2): [45.2 32.1 28.5 51.3 38.7]
    # Floor numbers: [0 1 2 3 4 5 6 7 8 9]
    # Column positions (m): [ 0  6 12 18 24 30]
    # Facade divisions (m): [ 0.  6. 12. 18. 24. 30.]
    # Initial occupancy: [0. 0. 0. 0. 0.]
    # Ceiling heights (m): [3.2 3.2 3.2 3.2 3.2 3.2 3.2 3.2]
    # 
    # Construction grid (0.1m):
    # Raw: [3.142 4.678 5.823 2.456]
    # Snapped: [3.1 4.7 5.8 2.5]
}{Different ways to create arrays for different needs}

\hh{Two-dimensional arrays: Working with building grids}

So far we've worked with one-dimensional arrays (like a single row or column in Excel). But buildings have floors and rooms---we need two-dimensional arrays (like a full spreadsheet):

\code{c14_2d_arrays_intro}{python}{
    import numpy as np
    
    # A 3-floor building with 4 apartments per floor
    # Each number represents the area in square meters
    apartment_areas = np.array([
        [45.0, 52.0, 48.0, 55.0],  # Floor 1
        [46.0, 53.0, 47.0, 54.0],  # Floor 2  
        [44.0, 51.0, 49.0, 56.0]   # Floor 3
    ])
    
    print("All apartment areas:")
    print(apartment_areas)
    
    # This is like a spreadsheet:
    # - 3 rows (floors)
    # - 4 columns (apartments per floor)
    # - 12 total values (3 x 4)
    
    # Understanding array properties
    print("\nArray properties:")
    print("Shape:", apartment_areas.shape)      # (3, 4) = 3 rows, 4 columns
    print("Total apartments:", apartment_areas.size)  # 12 apartments
    print("Dimensions:", apartment_areas.ndim)  # 2D array
    
    # Memory tip: axis=0 is down (vertical), axis=1 is across (horizontal)
    print("\nMemory tip: axis=0 is down, axis=1 is across")
    
    # Create a grid of zeros for a new building
    # 5 floors, 6 units per floor
    new_building = np.zeros((5, 6))
    print("\nNew building template shape:", new_building.shape)
    
    # Outcome:
    # All apartment areas:
    # [[45. 52. 48. 55.]
    #  [46. 53. 47. 54.]
    #  [44. 51. 49. 56.]]
    # 
    # Array properties:
    # Shape: (3, 4)
    # Total apartments: 12
    # Dimensions: 2
    # 
    # Memory tip: axis=0 is downward, axis=1 is left
    # 
    # New building template shape: (5, 6)
}{Think of 2D arrays as building floor plans}

\hh{Selecting data: Accessing specific rooms and floors}

Just like finding a specific cell in Excel (e.g., "B3"), you need to access specific elements in arrays:

\begin{tcolorbox}[
    title=Visual guide to array indexing,
    colframe=purple,
    colback = white,
    arc=2mm,
    breakable
]
\textbf{1D array indexing}

Let's understand how to select specific elements from a one-dimensional array. This is fundamental to all array operations.

\begin{lstlisting}[basicstyle=\ttfamily\footnotesize,breaklines=true,columns=fullflexible,keepspaces=true,xleftmargin=0pt]
Array:    [6.0,  7.5,  9.0,  7.5,  6.0]
Index:      0     1     2     3     4    (forward)
           -5    -4    -3    -2    -1    (backward)

spacings[0] = 6.0  (First element)
spacings[2] = 9.0  (Middle element)  
spacings[-1] = 6.0 (Last element)

Slicing: spacings[1:4]
         [6.0,  7.5,  9.0,  7.5,  6.0]
               start----end
           Returns: [7.5, 9.0, 7.5]
           (includes start, excludes end)
\end{lstlisting}

\textbf{Key indexing concepts:}
\begin{itemize}
\item \textbf{Forward indexing (0, 1, 2...):} Starts at 0, not 1. This is universal in programming. The first element is always at index 0.
\item \textbf{Backward indexing (-1, -2, -3...):} Negative indices count from the end. -1 is the last element, -2 is second-to-last, etc. Very useful when you don't know the array length.
\item \textbf{Slicing [start:end]:} Extracts a portion of the array. Crucially, it includes the start index but excludes the end index. So [1:4] gives you elements at positions 1, 2, and 3 (not 4).
\item \textbf{Why exclude the end?} This makes the math cleaner: [1:4] extracts exactly 4-1=3 elements.
\end{itemize}
\vspace{2em}
\textbf{2D array indexing}

Now let's see how to navigate a two-dimensional array, like finding a specific apartment in a building.

\begin{lstlisting}[basicstyle=\ttfamily\footnotesize,breaklines=true,columns=fullflexible,keepspaces=true,xleftmargin=0pt]
                Column index (0 to 3)
              0     1     2     3
         +------+------+------+------+
Row  0   |[0,0] |[0,1] |[0,2] |[0,3] |  Floor 0
Index    | 45.0 | 52.0 | 48.0 | 55.0 |
  |      +------+------+------+------+
  v  1   |[1,0] |[1,1] |[1,2] |[1,3] |  Floor 1
         | 46.0 | 53.0 | 47.0 | 54.0 |
         +------+------+------+------+
     2   |[2,0] |[2,1] |[2,2] |[2,3] |  Floor 2
         | 44.0 | 51.0 | 49.0 | 56.0 |
         +------+------+------+------+

apartments[1, 2] = 47.0
           ^   ^
         row  col
      (Floor 1, Unit 2)

Selecting entire row (floor):
apartments[1, :] gives [46.0, 53.0, 47.0, 54.0]
              ^
           (: means "all columns")

Selecting entire column (unit stack):
apartments[:, 2] gives [48.0, 47.0, 49.0]
           ^
        (: means "all rows")
\end{lstlisting}

\textbf{Understanding 2D indexing:}
\begin{itemize}
\item \textbf{[row, column] notation:} Always specify row first, then column. This matches mathematical matrix notation.
\item \textbf{The grid labels [0,0], [0,1], etc.:} Show the index pair for each cell. Think of it as coordinates.
\item \textbf{The colon (:) operator:} Means "all" or "everything". So [1, :] means "row 1, all columns" (entire floor 1).
\item \textbf{Real-world meaning:} apartments[1, 2] = "Floor 1, Unit 2" = 47.0 square meters.
\item \textbf{Common use cases:} Get all apartments on floor 2: apartments[2, :]. Get all corner units: apartments[:, 0] or apartments[:, -1].
\end{itemize}
\vspace{2em}
\textbf{Advanced slicing}

You can select rectangular regions of a 2D array, like selecting a block of cells in Excel.

\begin{lstlisting}[basicstyle=\ttfamily\footnotesize,breaklines=true,columns=fullflexible,keepspaces=true,xleftmargin=0pt]
Selecting a sub-region:
apartments[0:2, 1:3]  # Floors 0-1, Units 1-2

Original:                Selected region:
+----+----+----+----+   
| 45 | 52 | 48 | 55 |   +----+----+
+----+----+----+----+   | 52 | 48 |
| 46 | 53 | 47 | 54 |   +----+----+
+----+----+----+----+   | 53 | 47 |
| 44 | 51 | 49 | 56 |   +----+----+
+----+----+----+----+
      ^^^^
   Selected area
\end{lstlisting}
\vspace{2em}
\textbf{Sub-region selection explained:}
\begin{itemize}
\item \textbf{[0:2, 1:3] breakdown:} First part (0:2) selects rows 0 and 1. Second part (1:3) selects columns 1 and 2.
\item \textbf{The result:} A 2x2 grid containing just the middle apartments from the first two floors.
\item \textbf{Use cases:} Select a specific wing of the building, analyze only certain unit types, or focus on a problem area.
\item \textbf{Remember:} The end indices (2 and 3) are excluded, so we get rows 0-1 and columns 1-2.
\end{itemize}
\end{tcolorbox}

\code{c14_indexing_detailed}{python}{
    import numpy as np
    
    # One-dimensional example: 5 column spacings along a facade
    spacings = np.array([6.0, 7.5, 9.0, 7.5, 6.0])  # meters
    
    print("All spacings:", spacings)
    print("First spacing:", spacings[0])       # Remember: counting starts at 0!
    print("Last spacing:", spacings[-1])       # -1 means "last item"
    print("Middle three:", spacings[1:4])      # Items 1, 2, 3 (not 4!)
    
    # Two-dimensional example: apartment areas by floor
    apartments = np.array([
        [45.0, 52.0, 48.0, 55.0],  # Floor 0 (ground)
        [46.0, 53.0, 47.0, 54.0],  # Floor 1
        [44.0, 51.0, 49.0, 56.0]   # Floor 2
    ])
    
    # Selecting specific apartments
    print("\n--- Selecting specific data ---")
    print("Floor 0, Apt 1:", apartments[0, 1])  # 52.0 m^2
    print("Floor 2, Apt 3:", apartments[2, 3])  # 56.0 m^2
    
    # Selecting entire floors or apartment columns
    print("\nAll of Floor 1:", apartments[1, :])     # The : means "all"
    print("All Apartment 0's:", apartments[:, 0])   # First apt on each floor
    print("All Apartment 3's:", apartments[:, 3])   # Last apt on each floor
    
    # Selecting a range
    print("\nFloors 0-1, Apts 1-2:")
    print(apartments[0:2, 1:3])  # A 2x2 sub-grid
    
    # Outcome:
    # All spacings: [6.  7.5 9.  7.5 6. ]
    # First spacing: 6.0
    # Last spacing: 6.0
    # Middle three: [7.5 9.  7.5]
    # 
    # --- Selecting specific data ---
    # Floor 0, Apt 1: 52.0
    # Floor 2, Apt 3: 56.0
    # 
    # All of Floor 1: [46. 53. 47. 54.]
    # All Apartment 0's: [45. 46. 44.]
    # All Apartment 3's: [55. 54. 56.]
    # 
    # Floors 0-1, Apts 1-2:
    # [[52. 48.]
    #  [53. 47.]]
}{Like finding cells in a spreadsheet, but more powerful}

\hh{Filtering data: finding rooms that meet criteria}

One of NumPy's most powerful features is filtering---like using Excel's filter, but in code:

\begin{tcolorbox}[
    title=How filtering works,
    colframe=purple,
    colback = white,
    arc=2mm,
    breakable
]
\textbf{Step-by-step filtering process}

Let's understand how NumPy filters data by walking through a complete example of finding apartments that meet size criteria.

\begin{lstlisting}[basicstyle=\ttfamily\footnotesize,breaklines=true,columns=fullflexible,keepspaces=true,xleftmargin=0pt]
Step 1: Original array of apartment areas
areas = [28.0, 45.5, 52.0, 31.5, 48.0, 55.5, 42.0, 38.5]
         Unit0 Unit1 Unit2 Unit3 Unit4 Unit5 Unit6 Unit7

Step 2: Create boolean mask (areas > 45)
         [28.0, 45.5, 52.0, 31.5, 48.0, 55.5, 42.0, 38.5]
            |     |     |     |     |     |     |     |
         [False, True, True, False, True, True, False, False]
            X     /     /     X     /     /     X     X

Step 3: Apply mask to filter
areas[mask] --> [45.5, 52.0, 48.0, 55.5]
                Only values where mask is True
\end{lstlisting}

\textbf{Understanding the filtering process:}
\begin{itemize}
\item \textbf{Step 1:} We start with an array of apartment areas. Each has an index (Unit0, Unit1, etc.).
\item \textbf{Step 2:} The comparison (areas > 45) creates a "boolean mask"---an array of True/False values. Each True marks an apartment that meets our criteria.
\item \textbf{The symbols (X and /):} X marks rejected apartments (False), / marks accepted ones (True).
\item \textbf{Step 3:} When we apply the mask to the original array, NumPy extracts only the values where the mask is True.
\item \textbf{The result:} A new array containing only apartments larger than 45 square meters.
\end{itemize}
\vspace{2em}

\textbf{Multiple conditions}

Often you need to filter by multiple criteria, like finding medium-sized apartments.

\begin{lstlisting}[basicstyle=\ttfamily\footnotesize,breaklines=true,columns=fullflexible,keepspaces=true,xleftmargin=0pt]
Multiple conditions:
(areas >= 35) & (areas <= 50)
     |              |
     |              +-- Second condition: not too large
     +-- First condition: not too small
     
Step by step:
areas >= 35:    [F, T, T, F, T, T, T, T]  (not too small)
areas <= 50:    [T, T, F, T, T, F, T, T]  (not too large)
                           &               (both must be true)
Final mask:     [F, T, F, F, T, F, T, T]
Result:         [45.5, 48.0, 42.0, 38.5]
\end{lstlisting}

\textbf{Combining conditions:}
\begin{itemize}
\item \textbf{The \& operator:} Means "AND"---both conditions must be true. Use this for ranges.
\item \textbf{The | operator:} Means "OR"---either condition can be true. Use for alternatives.
\item \textbf{Parentheses are crucial:} Always put parentheses around each condition before combining.
\item \textbf{Step-by-step evaluation:} NumPy evaluates each condition separately, then combines the results.
\item \textbf{Common mistake:} Writing "35 <= areas <= 50" doesn't work in NumPy. You must split it into two conditions.
\end{itemize}
\vspace{2em}

\textbf{Visual example: finding studios}

Let's see how filtering works with categorical data, like apartment types.

\begin{lstlisting}[basicstyle=\ttfamily\footnotesize,breaklines=true,columns=fullflexible,keepspaces=true,xleftmargin=0pt]
unit_types = ['studio', 'one-bed', 'studio', 'two-bed', 
              'one-bed', 'studio', 'one-bed', 'two-bed']
areas =      [28.0,     45.5,     52.0,     31.5,
              48.0,     55.5,     42.0,     38.5]

Finding studios:
+--------+--------+--------+--------+--------+--------+--------+--------+
|studio  |one-bed |studio  |two-bed |one-bed |studio  |one-bed |two-bed |
| 28.0   | 45.5   | 52.0   | 31.5   | 48.0   | 55.5   | 42.0   | 38.5   |
+--------+--------+--------+--------+--------+--------+--------+--------+
    /        X        /        X        X        /        X        X
    
Result: studio_areas = [28.0, 52.0, 55.5]
\end{lstlisting}

\textbf{Filtering by category:}
\begin{itemize}
\item \textbf{Two parallel arrays:} One contains categories (unit types), the other contains numeric data (areas).
\item \textbf{The equality test:} unit\_types == 'studio' creates a mask marking which units are studios.
\item \textbf{Apply to areas:} We use this mask on the areas array to get only studio areas.
\item \textbf{Real-world use:} Find all corner units, all units facing north, all units on specific floors, etc.
\item \textbf{Combining with size:} You could further filter to find "large studios" by combining type and area conditions.
\end{itemize}
\end{tcolorbox}

\code{c14_filtering_detailed}{python}{
    import numpy as np
    
    # Room areas in square meters
    areas = np.array([28.0, 45.5, 52.0, 31.5, 48.0, 55.5, 42.0, 38.5])
    
    # STEP 1: Create a "mask" (True/False array)
    # Which apartments are larger than 45 square meters?
    large_mask = areas > 45.0
    print("Is each apartment > 45m^2?")
    print(large_mask)  # Array of True/False values
    
    # STEP 2: Use the mask to filter
    large_apartments = areas[large_mask]
    print("\nLarge apartments only:", large_apartments)
    
    # You can do this in one line:
    print("Medium apartments (35-50m^2):", areas[(areas >= 35) & (areas <= 50)])
    
    # Finding specific values
    unit_types = np.array(['studio', 'one-bed', 'studio', 'two-bed', 
                           'one-bed', 'studio', 'one-bed', 'two-bed'])
    
    # Which apartments are studios?
    studio_mask = unit_types == 'studio'
    studio_areas = areas[studio_mask]
    print("\nStudio areas:", studio_areas)
    
    # Multiple criteria: Large studios
    large_studio_mask = (unit_types == 'studio') & (areas > 40)
    print("Large studios:", areas[large_studio_mask])
    
    # Finding positions (indices) of matching elements
    positions = np.where(areas > 50)
    print("\nPositions of apartments > 50m^2:", positions[0])
    
    # Outcome:
    # Is each apartment > 45m^2?
    # [False  True  True False  True  True False False]
    # 
    # Large apartments only: [45.5 52.  48.  55.5]
    # Medium apartments (35-50m^2): [45.5 38.5 48. 42. 38.5]
    # 
    # Studio areas: [28. 52. 55.5]
    # Large studios: [52. 55.5]
    # 
    # Positions of apartments > 50m^2: [2 5]
}{Filter arrays like filtering spreadsheets, but with code}

\hh{Warning: views vs copies (avoiding a common mistake)}

This is a subtle but important concept. When you select part of an array, NumPy sometimes gives you a "view" (a window into the original data) rather than a copy. Changing the view changes the original!

\begin{tcolorbox}[
    title=Views vs copies illustrated,
    colframe=purple,
    colback = white,
    arc=2mm,
    breakable
]
\textbf{The problem: views share memory}

This is one of the most confusing aspects of NumPy for beginners. Let's understand exactly what happens when you select part of an array.

\begin{lstlisting}[basicstyle=\ttfamily\footnotesize,breaklines=true,columns=fullflexible,keepspaces=true,xleftmargin=0pt]
Original array in memory:
+-----+-----+-----+-----+-----+
|100.0|120.0|110.0|130.0|125.0|
+-----+-----+-----+-----+-----+
   0     1     2     3     4

Creating a view (selection = original[1:4]):
+-----+-----+-----+-----+-----+
|100.0|120.0|110.0|130.0|125.0|
+-----+--+--+--+--+--+--+-----+
         |      |      |
      +-----+-----+-----+
      |120.0|110.0|130.0|  <-- selection (VIEW)
      +-----+-----+-----+
      
Modifying the view (selection[0] = 999):
      +-----+-----+-----+
      |999.0|110.0|130.0|  <-- selection changed
      +-----+-----+-----+
         | (shares memory)
+-----+-----+-----+-----+-----+
|100.0|999.0|110.0|130.0|125.0|  <-- ORIGINAL CHANGED TOO!
+-----+-----+-----+-----+-----+
\end{lstlisting}

\textbf{Understanding the view problem:}
\begin{itemize}
\item \textbf{What's happening:} When you slice an array (like original[1:4]), NumPy doesn't immediately create new data. Instead, it creates a "view" that points to the same memory.
\item \textbf{Why NumPy does this:} For performance. Copying data takes time and memory. If you're just reading data, a view is much faster.
\item \textbf{The danger:} If you modify the view thinking it's independent, you accidentally change the original data!
\item \textbf{Real-world scenario:} You select floor 2 data to adjust for a renovation, but accidentally modify the master building database.
\item \textbf{The vertical line (|):} Shows the memory connection. The view and original share the same underlying data storage.
\end{itemize}
\vspace{2em}

\textbf{The solution: make a copy}

When you need to modify selected data without affecting the original, explicitly create a copy.

\begin{lstlisting}[basicstyle=\ttfamily\footnotesize,breaklines=true,columns=fullflexible,keepspaces=true,xleftmargin=0pt]
Creating a copy (selection = original[1:4].copy()):

Original memory:           Copy's memory:
+-----+-----+-----+-----+-----+   +-----+-----+-----+
|100.0|120.0|110.0|130.0|125.0|   |120.0|110.0|130.0|
+-----+-----+-----+-----+-----+   +-----+-----+-----+
                                   Independent copy!
                                   Different memory location

Modifying the copy (selection[0] = 777):
+-----+-----+-----+-----+-----+   +-----+-----+-----+
|100.0|120.0|110.0|130.0|125.0|   |777.0|110.0|130.0|
+-----+-----+-----+-----+-----+   +-----+-----+-----+
Original unchanged! o              Only copy changed
\end{lstlisting}
\vspace{2em}

\textbf{Understanding copies:}
\begin{itemize}
\item \textbf{The .copy() method:} Explicitly creates new, independent data in a different memory location.
\item \textbf{Two separate blocks:} The diagram shows two distinct memory regions with no connection.
\item \textbf{Safety:} Changes to the copy never affect the original.
\item \textbf{The cost:} Uses more memory and takes slightly longer to create.
\item \textbf{When to use:} Whenever you plan to modify the selected data, or when you need to ensure independence.
\end{itemize}
\vspace{2em}

\textbf{Quick check: is it a view?}

Here's how to check if an array is a view or an independent copy:

\begin{lstlisting}[basicstyle=\ttfamily\footnotesize,breaklines=true,columns=fullflexible,keepspaces=true,xleftmargin=0pt]
# Check if array is a view
if arr.base is not None:
    print("This is a VIEW - be careful!")
else:
    print("This is a COPY - safe to modify")

# The logic:
# - Views have a 'base' (the original array they view)
# - Copies have no base (base is None)
# - Original arrays also have no base
\end{lstlisting}
\vspace{2em}

\textbf{Practical guidelines:}
\begin{itemize}
\item \textbf{Reading only?} Views are fine and faster.
\item \textbf{Modifying data?} Always use .copy() to be safe.
\item \textbf{Unsure?} Use .copy() - better safe than sorry.
\item \textbf{Common operations that create views:} Basic slicing (arr[1:5]), selecting rows/columns (arr[2, :]).
\item \textbf{Operations that create copies:} Fancy indexing (arr[[1,3,5]]), boolean indexing (arr[arr>50]).
\end{itemize}
\end{tcolorbox}

\code{c14_view_copy_detailed}{python}{
    import numpy as np
    
    # Original floor areas
    original = np.array([100.0, 120.0, 110.0, 130.0, 125.0])
    print("Original areas:", original)
    
    # DANGER: Creating a view (not a copy)
    selection = original[1:4]  # Select middle three
    print("Selection:", selection)
    
    # If we modify the selection...
    selection[0] = 999.0  # Change first element of selection
    print("\nAfter changing selection[0] to 999:")
    print("Selection:", selection)
    print("Original:", original)  # OOPS! Original changed too!
    
    # SOLUTION: Use .copy() when you need independence
    original2 = np.array([100.0, 120.0, 110.0, 130.0, 125.0])
    safe_selection = original2[1:4].copy()  # Make an independent copy
    safe_selection[0] = 777.0
    
    print("\nWith .copy() method:")
    print("Safe selection:", safe_selection)
    print("Original2:", original2)  # Original unchanged!
    
    # Rule of thumb: Use .copy() when you plan to modify selected data
    
    # Outcome:
    # Original areas: [100. 120. 110. 130. 125.]
    # Selection: [120. 110. 130.]
    # 
    # After changing selection[0] to 999:
    # Selection: [999. 110. 130.]
    # Original: [100. 999. 110. 130. 125.]
    # 
    # With .copy() method:
    # Safe selection: [777. 110. 130.]
    # Original2: [100. 120. 110. 130. 125.]
}{Understanding when NumPy shares data vs. copying it}

\hh{Reshaping arrays: Reorganizing your data}

Sometimes you need to reorganize data---like taking a list of 12 apartments and arranging them into a 3x4 floor plan:

\code{c14_reshape_detailed}{python}{
    import numpy as np
    
    # Scenario: 12 apartment areas in a single list
    all_apartments = np.array([45, 52, 48, 55,   # These will become Floor 0
                               46, 53, 47, 54,   # These will become Floor 1
                               44, 51, 49, 56])  # These will become Floor 2
    
    print("Original shape:", all_apartments.shape)  # (12,) = one dimension
    
    # Reshape into 3 floors x 4 apartments
    building = all_apartments.reshape(3, 4)
    print("\nReshaped into building:")
    print(building)
    print("New shape:", building.shape)
    
    # You can also reshape into different configurations
    # As long as total elements stays the same (3x4 = 2x6 = 12)
    two_wings = all_apartments.reshape(2, 6)  # 2 wings, 6 units each
    print("\nAs 2 wings x 6 units:")
    print(two_wings)
    
    # Flatten back to 1D
    flat_again = building.flatten()
    print("\nFlattened back:", flat_again)
    
    # Use -1 to let NumPy figure out one dimension
    # "Make it 4 columns, you figure out the rows"
    auto_rows = all_apartments.reshape(-1, 4)
    print("\nAuto-calculated rows:", auto_rows.shape)  # (3, 4)
    
    # Transpose: flip rows and columns
    # Like rotating your spreadsheet 90 degrees
    transposed = building.T
    print("\nTransposed (4x3 instead of 3x4):")
    print(transposed)
    
    # Outcome:
    # Original shape: (12,)
    # 
    # Reshaped into building:
    # [[45 52 48 55]
    #  [46 53 47 54]
    #  [44 51 49 56]]
    # New shape: (3, 4)
    # 
    # As 2 wings x 6 units:
    # [[45 52 48 55 46 53]
    #  [47 54 44 51 49 56]]
    # 
    # Flattened back: [45 52 48 55 46 53 47 54 44 51 49 56]
    # 
    # Auto-calculated rows: (3, 4)
    # 
    # Transposed (4x3 instead of 3x4):
    # [[45 46 44]
    #  [52 53 51]
    #  [48 47 49]
    #  [55 54 56]]
}{Reorganize data like rearranging a spreadsheet}

\hh{Broadcasting: applying operations to entire arrays}

Broadcasting is NumPy's way of applying the same operation to every element without writing loops. It's like applying a formula to an entire column in Excel:

\begin{tcolorbox}[
    title=Broadcasting visualized,
    colframe=purple,
    colback = white,
    arc=2mm,
    breakable
]
\textbf{Single value broadcasting}

The simplest form of broadcasting is applying a single operation to every element in an array.

\begin{lstlisting}[basicstyle=\ttfamily\footnotesize,breaklines=true,columns=fullflexible,keepspaces=true,xleftmargin=0pt]
Converting all areas from m^2 to ft^2:

areas_m2 = [45.0, 52.0, 48.0, 55.0, 38.0]
               x     x     x     x     x
            10.764 10.764 10.764 10.764 10.764
               |     |     |     |     |
areas_ft2 = [484.4, 559.7, 516.7, 592.0, 409.0]

In NumPy: areas_ft2 = areas_m2 * 10.764
(One operation applies to all elements!)
\end{lstlisting}

\textbf{Understanding single value broadcasting:}
\begin{itemize}
\item \textbf{The multiplication symbols (x):} Show that the same conversion factor (10.764) is applied to each element.
\item \textbf{No explicit loop:} In traditional programming, you'd write a loop. NumPy handles this internally, much more efficiently.
\item \textbf{The single line:} areas\_ft2 = areas\_m2 * 10.764 replaces what would be 5+ lines of loop code.
\item \textbf{Works with any operation:} Addition, subtraction, division, powers---all work the same way.
\item \textbf{Real-world uses:} Apply safety factors, convert units, adjust for inflation, scale drawings, etc.
\end{itemize}
\vspace{2em}

\textbf{Array-to-array broadcasting}

When you operate on two arrays of the same size, NumPy pairs up corresponding elements.

\begin{lstlisting}[basicstyle=\ttfamily\footnotesize,breaklines=true,columns=fullflexible,keepspaces=true,xleftmargin=0pt]
Calculating room areas from dimensions:

widths =  [3.0,  3.5,  4.0,  3.8]
            x     x     x     x    <-- Multiplication happens
lengths = [8.0,  7.5,  6.0,  8.0]     element by element
            |     |     |     |
areas =   [24.0, 26.25, 24.0, 30.4]

Element-wise multiplication:
Position 0: 3.0 x 8.0 = 24.0
Position 1: 3.5 x 7.5 = 26.25
Position 2: 4.0 x 6.0 = 24.0
Position 3: 3.8 x 8.0 = 30.4
\end{lstlisting}

\textbf{Element-wise operations explained:}
\begin{itemize}
\item \textbf{Pairing:} Elements at the same position are paired: widths[0] with lengths[0], widths[1] with lengths[1], etc.
\item \textbf{The requirement:} Arrays must be the same size, or one must be broadcastable to the other's shape.
\item \textbf{Not matrix multiplication:} This is element-wise multiplication, different from mathematical matrix multiplication.
\item \textbf{Common operations:} Calculate areas from dimensions, combine loads and factors, apply different rates to different zones.
\end{itemize}
\vspace{2em}

\textbf{2D broadcasting with different shapes}

This is where broadcasting becomes truly powerful---applying different factors to different parts of your data.

\begin{lstlisting}[basicstyle=\ttfamily\footnotesize,breaklines=true,columns=fullflexible,keepspaces=true,xleftmargin=0pt]
Applying floor-specific reduction factors:

Base areas (3x4):          Floor factors (3x1):
+----+----+----+----+      +------+
| 45 | 52 | 48 | 55 |      | 0.85 | --> Applied to entire row
+----+----+----+----+      +------+
| 45 | 52 | 48 | 55 |  x   | 0.87 | --> Applied to entire row
+----+----+----+----+      +------+
| 45 | 52 | 48 | 55 |      | 0.89 | --> Applied to entire row
+----+----+----+----+      +------+

Result (net areas):
+------+------+------+------+
|38.25 |44.20 |40.80 |46.75 |  Floor 0 x 0.85
+------+------+------+------+
|39.15 |45.24 |41.76 |47.85 |  Floor 1 x 0.87
+------+------+------+------+
|40.05 |46.28 |42.72 |48.95 |  Floor 2 x 0.89
+------+------+------+------+

The (3x1) array "broadcasts" across each row!
\end{lstlisting}

\textbf{Shape broadcasting explained:}
\begin{itemize}
\item \textbf{The magic:} A (3x1) array (one value per floor) automatically expands to match the (3x4) array.
\item \textbf{How it works:} Each floor factor is applied to all apartments on that floor.
\item \textbf{The arrows:} Show that 0.85 is applied to all of floor 0, 0.87 to all of floor 1, etc.
\item \textbf{Real-world meaning:} Different floors might have different efficiency ratios, view premiums, or structural factors.
\item \textbf{The power:} One operation replaces nested loops that would be needed in traditional programming.
\end{itemize}
\vspace{2em}

\textbf{Broadcasting rules}

Not all array shapes can be broadcast together. Here are the rules:

\begin{lstlisting}[basicstyle=\ttfamily\footnotesize,breaklines=true,columns=fullflexible,keepspaces=true,xleftmargin=0pt]
Rule: Dimensions must be compatible
O (3, 4) with scalar     --> scalar expands to (3, 4)
O (3, 4) with (3, 1)     --> (3, 1) expands to (3, 4)
O (3, 4) with (1, 4)     --> (1, 4) expands to (3, 4)
X (3, 4) with (3, 3)     --> ERROR! Incompatible shapes

Compatible means:
- Dimensions are equal, OR
- One dimension is 1 (can be stretched), OR
- One array has fewer dimensions (can add dimensions of size 1)
\end{lstlisting}

\textbf{Understanding compatibility:}
\begin{itemize}
\item \textbf{Equal dimensions:} (3,4) and (3,4) obviously work together.
\item \textbf{Size 1 dimensions:} Can be "stretched" to match. (3,1) stretches to (3,4) by repeating.
\item \textbf{Scalars:} Always compatible---they broadcast to any shape.
\item \textbf{The error case:} (3,4) and (3,3) can't work because 4 $\neq$ 3 and neither is 1.
\item \textbf{Debugging tip:} When you get a broadcasting error, print the shapes of your arrays to understand the mismatch.
\end{itemize}
\end{tcolorbox}

\hh{Warning: views vs copies (avoiding a common mistake)}

This is a subtle but important concept. When you select part of an array, NumPy sometimes gives you a "view" (a window into the original data) rather than a copy. Changing the view changes the original!

\begin{tcolorbox}[
    title=Views vs copies illustrated,
    colframe=purple,
    colback = white,
    arc=2mm,
    breakable
]
\textbf{The problem: views share memory}

This is one of the most confusing aspects of NumPy for beginners. Let's understand exactly what happens when you select part of an array.

\begin{lstlisting}[basicstyle=\ttfamily\footnotesize,breaklines=true,columns=fullflexible,keepspaces=true,xleftmargin=0pt]
Original array in memory:
+-----+-----+-----+-----+-----+
|100.0|120.0|110.0|130.0|125.0|
+-----+-----+-----+-----+-----+
   0     1     2     3     4

Creating a view (selection = original[1:4]):
+-----+-----+-----+-----+-----+
|100.0|120.0|110.0|130.0|125.0|
+-----+--+--+--+--+--+--+-----+
         |      |      |
      +-----+-----+-----+
      |120.0|110.0|130.0|  <-- selection (VIEW)
      +-----+-----+-----+
      
Modifying the view (selection[0] = 999):
      +-----+-----+-----+
      |999.0|110.0|130.0|  <-- selection changed
      +-----+-----+-----+
         | (shares memory)
+-----+-----+-----+-----+-----+
|100.0|999.0|110.0|130.0|125.0|  <-- ORIGINAL CHANGED TOO!
+-----+-----+-----+-----+-----+
\end{lstlisting}

\textbf{Understanding the view problem:}
\begin{itemize}
\item \textbf{What's happening:} When you slice an array (like original[1:4]), NumPy doesn't immediately create new data. Instead, it creates a "view" that points to the same memory.
\item \textbf{Why NumPy does this:} For performance. Copying data takes time and memory. If you're just reading data, a view is much faster.
\item \textbf{The danger:} If you modify the view thinking it's independent, you accidentally change the original data!
\item \textbf{Real-world scenario:} You select floor 2 data to adjust for a renovation, but accidentally modify the master building database.
\item \textbf{The vertical line (|):} Shows the memory connection. The view and original share the same underlying data storage.
\end{itemize}
\vspace{2em}

\textbf{The solution: make a copy}

When you need to modify selected data without affecting the original, explicitly create a copy.

\begin{lstlisting}[basicstyle=\ttfamily\footnotesize,breaklines=true,columns=fullflexible,keepspaces=true,xleftmargin=0pt]
Creating a copy (selection = original[1:4].copy()):

Original memory:           Copy's memory:
+-----+-----+-----+-----+-----+   +-----+-----+-----+
|100.0|120.0|110.0|130.0|125.0|   |120.0|110.0|130.0|
+-----+-----+-----+-----+-----+   +-----+-----+-----+
                                   Independent copy!
                                   Different memory location

Modifying the copy (selection[0] = 777):
+-----+-----+-----+-----+-----+   +-----+-----+-----+
|100.0|120.0|110.0|130.0|125.0|   |777.0|110.0|130.0|
+-----+-----+-----+-----+-----+   +-----+-----+-----+
Original unchanged! O              Only copy changed
\end{lstlisting}

\textbf{Understanding copies:}
\begin{itemize}
\item \textbf{The .copy() method:} Explicitly creates new, independent data in a different memory location.
\item \textbf{Two separate blocks:} The diagram shows two distinct memory regions with no connection.
\item \textbf{Safety:} Changes to the copy never affect the original.
\item \textbf{The cost:} Uses more memory and takes slightly longer to create.
\item \textbf{When to use:} Whenever you plan to modify the selected data, or when you need to ensure independence.
\end{itemize}

\textbf{Quick check: is it a view?}

Here's how to check if an array is a view or an independent copy:

\begin{lstlisting}[basicstyle=\ttfamily\footnotesize,breaklines=true,columns=fullflexible,keepspaces=true,xleftmargin=0pt]
# Check if array is a view
if arr.base is not None:
    print("This is a VIEW - be careful!")
else:
    print("This is a COPY - safe to modify")

# The logic:
# - Views have a 'base' (the original array they view)
# - Copies have no base (base is None)
# - Original arrays also have no base
\end{lstlisting}
\vspace{2em}

\textbf{Practical guidelines:}
\begin{itemize}
\item \textbf{Reading only?} Views are fine and faster.
\item \textbf{Modifying data?} Always use .copy() to be safe.
\item \textbf{Unsure?} Use .copy() - better safe than sorry.
\item \textbf{Common operations that create views:} Basic slicing (arr[1:5]), selecting rows/columns (arr[2, :]).
\item \textbf{Operations that create copies:} Fancy indexing (arr[[1,3,5]]), boolean indexing (arr[arr>50]).
\end{itemize}
\end{tcolorbox}

\code{c14_view_copy_detailed}{python}{
    import numpy as np
    
    # Original floor areas
    original = np.array([100.0, 120.0, 110.0, 130.0, 125.0])
    print("Original areas:", original)
    
    # DANGER: Creating a view (not a copy)
    selection = original[1:4]  # Select middle three
    print("Selection:", selection)
    
    # If we modify the selection...
    selection[0] = 999.0  # Change first element of selection
    print("\nAfter changing selection[0] to 999:")
    print("Selection:", selection)
    print("Original:", original)  # OOPS! Original changed too!
    
    # SOLUTION: Use .copy() when you need independence
    original2 = np.array([100.0, 120.0, 110.0, 130.0, 125.0])
    safe_selection = original2[1:4].copy()  # Make an independent copy
    safe_selection[0] = 777.0
    
    print("\nWith .copy() method:")
    print("Safe selection:", safe_selection)
    print("Original2:", original2)  # Original unchanged!
    
    # Rule of thumb: Use .copy() when you plan to modify selected data
    
    # Outcome:
    # Original areas: [100. 120. 110. 130. 125.]
    # Selection: [120. 110. 130.]
    # 
    # After changing selection[0] to 999:
    # Selection: [999. 110. 130.]
    # Original: [100. 999. 110. 130. 125.]
    # 
    # With .copy() method:
    # Safe selection: [777. 110. 130.]
    # Original2: [100. 120. 110. 130. 125.]
}{Understanding when NumPy shares data vs. copying it}

\hh{Common mistakes architecture students make}

\begin{tcolorbox}[
    title=Avoid These Common Pitfalls,
    colframe=purple,
    colback = white,
    arc=2mm
]
\textbf{Mistake 1: Forgetting to label units}
\begin{lstlisting}[basicstyle=\ttfamily\footnotesize,breaklines=true,columns=fullflexible,keepspaces=true,xleftmargin=0pt]
# BAD: What units are these?
areas = np.array([45, 52, 48])

# GOOD: Clear about units
areas_m2 = np.array([45, 52, 48])  # square meters
areas_ft2 = areas_m2 * 10.764      # converted to square feet

\end{lstlisting}

\textbf{Mistake 2: Wrong axis for aggregation}
\begin{lstlisting}[basicstyle=\ttfamily\footnotesize,breaklines=true,columns=fullflexible,keepspaces=true,xleftmargin=0pt]
# CONFUSING: Which axis is which?
building.sum(axis=0)  # Sums vertically (by unit position)
building.sum(axis=1)  # Sums horizontally (by floor) - usually what you want

# MEMORY TIP: axis=0 is down, axis=1 is across
floor_totals = building.sum(axis=1)  # Sum across each floor

\end{lstlisting}

\textbf{Mistake 3: Not considering construction constraints}
\begin{lstlisting}[basicstyle=\ttfamily\footnotesize,breaklines=true,columns=fullflexible,keepspaces=true,xleftmargin=0pt]
# BAD: Unrealistic precision
wall_length = 3.142857 

# GOOD: Snap to construction grid
wall_length = np.round(3.142857, 1)  # 3.1m (0.1m grid)
wall_length = np.round(3.142857 * 2) / 2  # 3.0m (0.5m grid)

\end{lstlisting}

\textbf{Mistake 4: Modifying views instead of copies}
\begin{lstlisting}[basicstyle=\ttfamily\footnotesize,breaklines=true,columns=fullflexible,keepspaces=true,xleftmargin=0pt]
# DANGEROUS: This modifies the original!
selection = original_data[1:4]
selection[0] = 999  # Changes original_data too!

# SAFE: Make a copy first
selection = original_data[1:4].copy()
selection[0] = 999  # Original unchanged

\end{lstlisting}
\end{tcolorbox}

\hh{Aggregation: Calculating totals and statistics}

NumPy makes it easy to calculate statistics for your entire dataset or by floor/wing:

\code{c14_aggregation_detailed}{python}{
    import numpy as np
    
    # Building with 4 floors, 5 units per floor
    apartment_areas = np.array([
        [45.2, 52.1, 48.3, 55.0, 42.5],  # Ground floor
        [46.0, 53.2, 47.5, 54.3, 43.0],  # Floor 1
        [44.8, 51.5, 49.0, 56.2, 41.8],  # Floor 2
        [45.5, 52.8, 48.8, 55.5, 42.0]   # Floor 3
    ])
    
    # Basic statistics for the entire building
    print("=== BUILDING STATISTICS ===")
    print(f"Total area: {apartment_areas.sum():.1f} m^2")
    print(f"Average apartment: {apartment_areas.mean():.1f} m^2")
    print(f"Smallest apartment: {apartment_areas.min():.1f} m^2")
    print(f"Largest apartment: {apartment_areas.max():.1f} m^2")
    print(f"Standard deviation: {apartment_areas.std():.1f} m^2")
    
    # Statistics by floor (axis=1 means "across the columns")
    print("\n=== BY FLOOR (axis=1 across) ===")
    floor_totals = apartment_areas.sum(axis=1)
    floor_averages = apartment_areas.mean(axis=1)
    print("Total area per floor:", floor_totals)
    print("Average apartment per floor:", floor_averages)
    
    # Statistics by apartment position (axis=0 means "down the rows")
    print("\n=== BY APARTMENT POSITION (axis=0 down) ===")
    position_averages = apartment_areas.mean(axis=0)
    print("Average by position (0-4):", position_averages)
    print("(e.g., corner units vs. middle units)")
    
    # Finding specific values
    print("\n=== FINDING APARTMENTS ===")
    # Which apartments are exactly 45.5 m^2?
    matches = np.where(apartment_areas == 45.5)
    print(f"Apartment at floor {matches[0][0]}, unit {matches[1][0]} is 45.5 m^2")
    
    # Top 3 largest apartments
    flat_areas = apartment_areas.flatten()
    sorted_areas = np.sort(flat_areas)
    top_3 = sorted_areas[-3:]  # Last 3 elements
    print(f"Top 3 largest: {top_3} m^2")
    
    # Outcome:
    # === BUILDING STATISTICS ===
    # Total area: 968.2 m^2
    # Average apartment: 48.4 m^2
    # Smallest apartment: 41.8 m^2
    # Largest apartment: 56.2 m^2
    # Standard deviation: 5.1 m^2
    # 
    # === BY FLOOR (axis=1 across) ===
    # Total area per floor: [243.1 244.  243.3 244.6]
    # Average apartment per floor: [48.62 48.8  48.66 48.92]
    # 
    # === BY APARTMENT POSITION (axis=0 down) ===
    # Average by position (0-4): [45.375 52.4   48.4   55.25  42.325]
    # (e.g., corner units vs. middle units)
    # 
    # === FINDING APARTMENTS ===
    # Apartment at floor 3, unit 0 is 45.5 m^2
    # Top 3 largest: [55.  55.5 56.2] m^2
}{Calculate statistics for buildings and floors}

% ====================================================================
% PART II -- MATHEMATICAL FUNCTIONS FOR ARCHITECTURE
% ====================================================================

\hh{Mathematical operations for design}

NumPy includes many mathematical functions that apply to entire arrays at once. Here are the most useful ones for architectural work:

\code{c14_math_functions}{python}{
    import numpy as np
    
    # Basic arithmetic on arrays
    spans = np.array([6.0, 7.5, 9.0, 7.5, 6.0])  # Beam spans in meters
    
    # All operations apply to every element
    print("Original spans:", spans)
    print("Double spans:", spans * 2)
    print("Half spans:", spans / 2)
    print("Spans + 1m:", spans + 1)
    print("Spans squared:", spans ** 2)
    
    # Combining arrays
    loads = np.array([10.0, 12.0, 15.0, 12.0, 10.0])  # kN/m
    moments = (loads * spans**2) / 8  # Simple beam formula
    print("\nMax moments:", moments, "kN*m")
    
    # Rounding for practical dimensions
    raw_dimensions = np.array([3.142, 4.678, 5.823, 2.456])
    print("\nRaw dims:", raw_dimensions)
    print("Rounded:", np.round(raw_dimensions, 1))      # To 1 decimal
    print("Rounded to 0.5:", np.round(raw_dimensions * 2) / 2)  # To nearest 0.5
    print("Ceiling:", np.ceil(raw_dimensions))          # Round up
    print("Floor:", np.floor(raw_dimensions))            # Round down
    
    # Outcome:
    # Original spans: [6.  7.5 9.  7.5 6. ]
    # Double spans: [12. 15. 18. 15. 12.]
    # Half spans: [3.   3.75 4.5  3.75 3.  ]
    # Spans + 1m: [ 7.   8.5 10.   8.5  7. ]
    # Spans squared: [36.   56.25 81.   56.25 36.  ]
    # 
    # Max moments: [45.    84.375 151.875 84.375 45.   ] kN-m
    # 
    # Raw dims: [3.142 4.678 5.823 2.456]
    # Rounded: [3.1 4.7 5.8 2.5]
    # Rounded to 0.5: [3.  4.5 6.  2.5]
    # Ceiling: [4. 5. 6. 3.]
    # Floor: [3. 4. 5. 2.]
}{Common mathematical operations for design}

\hh{Working with angles and trigonometry}

For sun angles, ramp slopes, and roof pitches, you'll need trigonometric functions:

\begin{tcolorbox}[
    title=Trigonometry in engineering and architecture,
    colframe=purple,
    colback = white,
    arc=2mm
]
\textbf{Ramp Slope Calculation}
\begin{lstlisting}[basicstyle=\ttfamily\footnotesize,breaklines=true,columns=fullflexible,keepspaces=true,xleftmargin=0pt]
Side View of Ramp:
                          
     Rise |               angle = 10 deg
          |            /
          |         /
          |      /     slope = tan(angle) x 100%
          |   /              = tan(10 deg) x 100%
          |/ angle           = 17.6%
    ------+----------
      Horizontal Run
        (10m)

For 10m run at different angles:
Angle:  5 deg  7.5 deg  10 deg  12.5 deg
Rise:   0.87m  1.32m    1.76m   2.22m
Slope:  8.7%   13.2%    17.6%   22.2%

\end{lstlisting}

\textbf{Sun Shading Calculation}
\begin{lstlisting}[basicstyle=\ttfamily\footnotesize,breaklines=true,columns=fullflexible,keepspaces=true,xleftmargin=0pt]
Section Through Window:

Overhang Depth Calculation:
                    
    |------d------| Overhang depth needed
    =============== (d = h / tan(sun_angle))
          \
           \ sun ray
            \
    +--------\----+
    |         \   |
    | Window  \  | h = 2.5m
    |          \ |
    |       angle\|
    +-------------+
    
Sun Angle vs Overhang Needed:
15 deg gives 9.33m (early morning/late afternoon)
30 deg gives 4.33m 
45 deg gives 2.50m (mid-morning/mid-afternoon)
60 deg gives 1.44m 
75 deg gives 0.67m (noon in summer)

\end{lstlisting}

\textbf{Degrees vs Radians}
\begin{lstlisting}[basicstyle=\ttfamily\footnotesize,breaklines=true,columns=fullflexible,keepspaces=true,xleftmargin=0pt]
Important: NumPy uses RADIANS for trig functions!

Degrees:  0    30    45    60    90     180    360
          |    |     |     |     |      |      |
Radians:  0   pi/6  pi/4  pi/3  pi/2    pi    2*pi
         0.0  0.52  0.79  1.05  1.57   3.14   6.28

Convert: rads = np.deg2rad(degrees)
         degs = np.rad2deg(radians)

\end{lstlisting}
\end{tcolorbox}

\code{c14_trigonometry}{python}{
    import numpy as np
    
    # Ramp slopes in degrees
    ramp_angles_deg = np.array([5.0, 7.5, 10.0, 12.5])
    
    # Convert to radians (NumPy trig functions use radians)
    ramp_angles_rad = np.deg2rad(ramp_angles_deg)
    
    # Calculate rise for 10m horizontal run
    horizontal_run = 10.0  # meters
    rises = horizontal_run * np.tan(ramp_angles_rad)
    slopes_percent = np.tan(ramp_angles_rad) * 100
    
    print("Ramp angles (degrees):", ramp_angles_deg)
    print("Rise over 10m run:", np.round(rises, 2), "m")
    print("Slope percentage:", np.round(slopes_percent, 1), "%")
    
    # Sun angles for shading calculation
    # Solar altitude angles at different times
    sun_altitudes_deg = np.array([15, 30, 45, 60, 75])  # degrees
    sun_altitudes_rad = np.deg2rad(sun_altitudes_deg)
    
    # Overhang depth needed for 2.5m window height
    window_height = 2.5  # meters
    overhang_depths = window_height / np.tan(sun_altitudes_rad)
    
    print("\nSun altitude:", sun_altitudes_deg, "degrees")
    print("Overhang depth needed:", np.round(overhang_depths, 2), "m")
    
    # Outcome:
    # Ramp angles (degrees): [ 5.   7.5 10.  12.5]
    # Rise over 10m run: [0.87 1.32 1.76 2.22] m
    # Slope percentage: [ 8.7 13.2 17.6 22.2] %
    # 
    # Sun altitude: [15 30 45 60 75] degrees
    # Overhang depth needed: [9.33 4.33 2.5  1.44 0.67] m
}{Trigonometry for slopes and sun angles}

\hh{Calculating differences and cumulative values}

For floor-to-floor heights, accumulated areas, and progressive dimensions:

\code{c14_diff_cumsum}{python}{
    import numpy as np
    
    # Floor elevations in meters
    floor_elevations = np.array([0.0, 3.3, 6.6, 10.0, 13.4, 16.8])
    
    # Calculate floor-to-floor heights
    floor_heights = np.diff(floor_elevations)
    print("Floor elevations:", floor_elevations)
    print("Floor-to-floor heights:", floor_heights)
    print("Average floor height:", np.mean(floor_heights), "m")
    
    # Progressive area calculation
    # Areas of spaces as you move through building
    room_areas = np.array([25.0, 30.0, 45.0, 28.0, 35.0, 40.0])
    cumulative_area = np.cumsum(room_areas)
    
    print("\nIndividual room areas:", room_areas)
    print("Cumulative area:", cumulative_area)
    print("(Total area after each room)")
    
    # Grid spacing accumulation
    column_spacings = np.array([6.0, 7.5, 7.5, 7.5, 6.0])  # meters
    column_positions = np.insert(np.cumsum(column_spacings), 0, 0)
    
    print("\nColumn spacings:", column_spacings)
    print("Column positions from start:", column_positions)
    
    # Outcome:
    # Floor elevations: [ 0.   3.3  6.6 10.  13.4 16.8]
    # Floor-to-floor heights: [3.3 3.3 3.4 3.4 3.4]
    # Average floor height: 3.36 m
    # 
    # Individual room areas: [25. 30. 45. 28. 35. 40.]
    # Cumulative area: [ 25.  55. 100. 128. 163. 203.]
    # (Total area after each room)
    # 
    # Column spacings: [6.  7.5 7.5 7.5 6. ]
    # Column positions from start: [ 0.   6.  13.5 21.  28.5 34.5]
}{Track differences and running totals}

% ====================================================================
% PART III -- PRACTICAL ARCHITECTURAL EXAMPLES
% ====================================================================

\hh{Mini-Lab 1: Apartment Building Analysis}

Let's combine everything we've learned to analyze a small apartment building:

\code{c14_example_building}{python}{
    import numpy as np
    
    # 5-story building, 4 apartments per floor
    # Each value is the apartment area in m^2
    building = np.array([
        [55.0, 72.0, 68.0, 85.0],  # Ground floor
        [56.0, 73.0, 67.0, 84.0],  # Floor 1
        [54.0, 71.0, 69.0, 86.0],  # Floor 2
        [55.5, 72.5, 68.5, 85.5],  # Floor 3
        [56.5, 73.5, 67.5, 84.5]   # Floor 4 (top)
    ])
    
    print("=== BUILDING ANALYSIS ===\n")
    
    # Basic statistics
    print(f"Total building area: {building.sum():.0f} m^2")
    print(f"Number of apartments: {building.size}")
    print(f"Average apartment size: {building.mean():.1f} m^2")
    print(f"Smallest apartment: {building.min():.0f} m^2")
    print(f"Largest apartment: {building.max():.0f} m^2")
    
    # Analysis by floor
    floor_totals = building.sum(axis=1)
    print(f"\nFloor totals: {floor_totals} m^2")
    print(f"Most area on floor: {np.argmax(floor_totals)}")
    
    # Analysis by apartment stack (vertical)
    stack_averages = building.mean(axis=0)
    print(f"\nAverage by stack position: {stack_averages}")
    print("(Shows if corner/middle units differ)")
    
    # Find apartments meeting criteria
    print("\n=== APARTMENT SEARCH ===")
    
    # Studios (< 60 m^2)
    studio_mask = building < 60
    num_studios = studio_mask.sum()
    print(f"Studios (<60m^2): {num_studios} units")
    
    # Large apartments (> 80 m^2)
    large_mask = building > 80
    large_apts = building[large_mask]
    print(f"Large apartments (>80m^2): {large_apts}")
    
    # Which floors have apartments between 70-75 m^2?
    mid_range = (building >= 70) & (building <= 75)
    floors_with_mid = np.any(mid_range, axis=1)
    floor_numbers = np.where(floors_with_mid)[0]
    print(f"Floors with 70-75m^2 apartments: {floor_numbers}")
    
    # Price calculation
    price_per_m2 = 3500  # $/m^2
    apartment_prices = building * price_per_m2
    
    print(f"\n=== PRICING ===")
    print(f"Total building value: ${apartment_prices.sum():,.0f}")
    print(f"Average apartment price: ${apartment_prices.mean():,.0f}")
    print(f"Floor 0 prices: ${apartment_prices[0]}")
    
    # Outcome:
    # === BUILDING ANALYSIS ===
    # 
    # Total building area: 1400 m^2
    # Number of apartments: 20
    # Average apartment size: 70.0 m^2
    # Smallest apartment: 54 m^2
    # Largest apartment: 86 m^2
    # 
    # Floor totals: [280. 280. 280. 282. 282.] m^2
    # Most area on floor: 3
    # 
    # Average by stack position: [55.4 72.4 68.  85. ]
    # (Shows if corner/middle units differ)
    # 
    # === APARTMENT SEARCH ===
    # Studios (<60m^2): 8 units
    # Large apartments (>80m^2): [85. 84. 86. 85.5 84.5]
    # Floors with 70-75m^2 apartments: [0 1 2 3 4]
    # 
    # === PRICING ===
    # Total building value: $4,900,000
    # Average apartment price: $245,000
    # Floor 0 prices: $[192500. 252000. 238000. 297500.]
}{Complete building analysis workflow}

\hh{Facade Panel Generation}

Creating a grid of panels for a parametric facade:

\begin{tcolorbox}[
    title=Facade Grid Generation (Visual),
    colframe=purple,
    colback = white,
    arc=2mm
]
\textbf{Creating the Panel Grid}
\begin{lstlisting}[basicstyle=\ttfamily\footnotesize,breaklines=true,columns=fullflexible,keepspaces=true,xleftmargin=0pt]
Building Facade (30m \times 12m):
    
    0    1.5   3.0   4.5   ...   27.0  28.5  30.0m
12m +----+----+----+----+----+----+----+----+----+
    |    |    |    |    |    |    |    |    |    |
10  +----+----+----+----+----+----+----+----+----+
    |    |    |    |    |    |    |    |    |    |
8   +----+----+----+----+----+----+----+----+----+
    |    |    |    | x  |    |    |    |    |    |
6   +----+----+----+----+----+----+----+----+----+
    |    |    |    |    |    |    |    |    |    |
4   +----+----+----+----+----+----+----+----+----+
    |    | x  |    |    |    |    |    |    |    |
2   +----+----+----+----+----+----+----+----+----+
    |    |    |    |    |    |    |    |    |    |
0m  +----+----+----+----+----+----+----+----+----+
    
    Panel size: 1.5m \times 2.0m
    Grid: 20 panels \times 6 panels = 120 total
    x = Panel center points

Panel Centers Calculation:
X: start at panel_width/2, step by panel_width
   [0.75, 2.25, 3.75, 5.25, ... 28.75]
   
Y: start at panel_height/2, step by panel_height
   [1.0, 3.0, 5.0, 7.0, 9.0, 11.0]

\end{lstlisting}

\textbf{Height-Based Rotation Pattern}
\begin{lstlisting}[basicstyle=\ttfamily\footnotesize,breaklines=true,columns=fullflexible,keepspaces=true,xleftmargin=0pt]
Side View - Panel Rotation by Height:

12m +-+ <- 30$^\circ$ rotation (maximum shading)
    |/|
10  +-+
    |/| <- Rotation increases with height
8   +-+    (better sun protection higher up)
    |/|
6   +-+
    |/|
4   +-+
    |/|
2   +-+
    ||| <- 0$^\circ$ rotation (minimum shading)
0m  +-+
    
Formula: rotation = (height/total_height) \times 30$^\circ$
         + random(-5$^\circ$, +5$^\circ$) for variation

Opening % = 100 - (rotation/45) \times 50

\end{lstlisting}
\end{tcolorbox}

\begin{tcolorbox}[
    title=Facade Grid Generation (Visual) - cont'd,
    colframe=purple,
    colback = white,
    arc=2mm
]
\textbf{Entrance Zone Selection}
\begin{lstlisting}[mathescape=true][basicstyle=\ttfamily\footnotesize,breaklines=true,columns=fullflexible,keepspaces=true,xleftmargin=0pt]
Finding panels in entrance zone (12m \leq x \leq 18m, y \leq 4m):

    12m      18m
    $\downarrow$         $\downarrow$
+----+----+----+----+----+----+----+----+----+
|    |    |    |    |    |    |    |    |    | 12m
+----+----+----+----+----+----+----+----+----+
|    |    |    |    |    |    |    |    |    | 10
+----+----+----+----+----+----+----+----+----+
|    |    |    |    |    |    |    |    |    | 8
+----+----+----+----+----+----+----+----+----+
|    |    |    |    |    |    |    |    |    | 6
+----+----+----+----+----+----+----+----+----+ \leftarrow 4m
|    |    | $\blacksquare$  | $\blacksquare$  | $\blacksquare$  | $\blacksquare$  |    |    |    | 2
+----+----+----+----+----+----+----+----+----+
|    |    | $\blacksquare$  | $\blacksquare$  | $\blacksquare$  | $\blacksquare$  |    |    |    | 0m
+----+----+----+----+----+----+----+----+----+
           +----------+
         Entrance Zone
         (8 panels marked with $\blacksquare$)

\end{lstlisting}
\end{tcolorbox}

\code{c14_example_facade}{python}{
    import numpy as np
    
    # Building facade dimensions
    width = 30.0   # meters
    height = 12.0  # meters
    
    # Panel dimensions
    panel_width = 1.5   # meters
    panel_height = 2.0  # meters
    
    # Generate panel center points
    # Start at half-panel distance from edge
    x_positions = np.arange(panel_width/2, width, panel_width)
    y_positions = np.arange(panel_height/2, height, panel_height)
    
    print(f"Panels in X direction: {len(x_positions)}")
    print(f"Panels in Y direction: {len(y_positions)}")
    print(f"Total panels: {len(x_positions) * len(y_positions)}")
    
    # Create grid of all panel centers
    X, Y = np.meshgrid(x_positions, y_positions)
    
    # Flatten to get list of (x, y) coordinates
    panel_centers = np.column_stack([X.flatten(), Y.flatten()])
    
    print(f"\nFirst 5 panel centers:")
    print(panel_centers[:5])
    
    # Add variation: rotate panels based on height
    # Higher panels get more rotation (for sun shading)
    heights = panel_centers[:, 1]  # Y coordinates
    
    # Rotation angle increases with height (0 deg at bottom, 30 deg at top)
    rotation_angles = (heights / height) * 30  # degrees
    
    # Add random variation (+/- 5 degrees)
    np.random.seed(42)  # For reproducibility
    random_variation = np.random.uniform(-5, 5, len(rotation_angles))
    final_rotations = rotation_angles + random_variation
    
    # Ensure rotations stay in range [0, 45] degrees
    final_rotations = np.clip(final_rotations, 0, 45)
    
    print(f"\nSample rotations for panels at different heights:")
    sample_indices = [0, 30, 60, 90, 119]  # Bottom to top
    for i in sample_indices:
        print(f"Panel {i:3d} at height {heights[i]:.1f}m: "
              f"{final_rotations[i]:.1f} deg")
    
    # Calculate opening percentage based on rotation
    # 0 deg = 100% open, 45 deg = 50% open
    opening_percentages = 100 - (final_rotations / 45) * 50
    
    print(f"\nAverage opening: {opening_percentages.mean():.1f}%")
    print(f"Min opening: {opening_percentages.min():.1f}%")
    print(f"Max opening: {opening_percentages.max():.1f}%")
    
    # Find panels in a specific zone (e.g., entrance area)
    entrance_zone = (panel_centers[:, 0] >= 12) & \
                   (panel_centers[:, 0] <= 18) & \
                   (panel_centers[:, 1] <= 4)
    
    entrance_panels = panel_centers[entrance_zone]
    print(f"\nPanels in entrance zone: {len(entrance_panels)}")
    print("Entrance panel positions:")
    print(entrance_panels)
    
    # Outcome:
    # Panels in X direction: 20
    # Panels in Y direction: 6
    # Total panels: 120
    # 
    # First 5 panel centers:
    # [[0.75 1.  ]
    #  [2.25 1.  ]
    #  [3.75 1.  ]
    #  [5.25 1.  ]
    #  [6.75 1.  ]]
    # 
    # Sample rotations for panels at different heights:
    # Panel   0 at height 1.0m: 2.1 deg
    # Panel  30 at height 3.0m: 5.5 deg
    # Panel  60 at height 5.0m: 13.2 deg
    # Panel  90 at height 7.0m: 19.5 deg
    # Panel 119 at height 11.0m: 28.7 deg
    # 
    # Average opening: 82.3%
    # Min opening: 68.1%
    # Max opening: 100.0%
    # 
    # Panels in entrance zone: 8
    # Entrance panel positions:
    # [[12.75  1.  ]
    #  [14.25  1.  ]
    #  [15.75  1.  ]
    #  [17.25  1.  ]
    #  [12.75  3.  ]
    #  [14.25  3.  ]
    #  [15.75  3.  ]
    #  [17.25  3.  ]]
}{Generate parametric facade with variations}

% ====================================================================
% PART IV -- INTEGRATION WITH OTHER TOOLS
% ====================================================================

\hh{Working with external data}

NumPy seamlessly connects with other tools you use in architectural practice:

\code{c14_external_data}{python}{
    import numpy as np
    
    # Reading from CSV files (common export from Excel/Revit)
    # Suppose you have a file 'room_schedule.csv' with room data
    # rooms = np.loadtxt('room_schedule.csv', delimiter=',', skiprows=1)
    
    # For this example, we'll simulate the data
    print("=== READING EXTERNAL DATA ===")
    
    # Simulated room data: [room_id, area, perimeter, height]
    room_data = np.array([
        [101, 45.5, 27.2, 3.0],
        [102, 38.2, 24.8, 3.0],
        [103, 52.1, 29.0, 3.5],
        [104, 41.7, 26.0, 3.0],
        [105, 48.3, 28.0, 3.5]
    ])
    
    # Extract columns
    room_ids = room_data[:, 0].astype(int)
    areas = room_data[:, 1]
    perimeters = room_data[:, 2]
    heights = room_data[:, 3]
    
    # Calculate volumes
    volumes = areas * heights
    
    # Calculate wall areas (for painting/finishes)
    wall_areas = perimeters * heights
    
    print(f"Room IDs: {room_ids}")
    print(f"Volumes: {volumes} m^3")
    print(f"Wall areas: {wall_areas} m^2")
    
    # Save results back to CSV
    results = np.column_stack([room_ids, areas, volumes, wall_areas])
    # np.savetxt('room_analysis.csv', results, delimiter=',',
    #            header='ID,Area,Volume,WallArea', comments='')
    
    print("\n=== PREPARING DATA FOR VISUALIZATION ===")
    
    # Normalize data for color mapping (0-1 range)
    normalized_areas = (areas - areas.min()) / (areas.max() - areas.min())
    print(f"Normalized areas for coloring: {normalized_areas}")
    
    # Create categories for room sizes
    small = areas < 40
    medium = (areas >= 40) & (areas < 50)
    large = areas >= 50
    
    categories = np.zeros_like(areas)
    categories[medium] = 1
    categories[large] = 2
    
    print(f"Size categories (0=small, 1=medium, 2=large): {categories.astype(int)}")
    
    # Outcome:
    # === READING EXTERNAL DATA ===
    # Room IDs: [101 102 103 104 105]
    # Volumes: [136.5 114.6 182.35 125.1 169.05] m^3
    # Wall areas: [81.6 74.4 101.5 78. 98.] m^2
    # 
    # === PREPARING DATA FOR VISUALIZATION ===
    # Normalized areas for coloring: [0.46428571 0. 1. 0.25 0.71428571]
    # Size categories (0=small, 1=medium, 2=large): [1 0 2 1 1]
}{Connect NumPy with Excel, Revit, and visualization tools}

% ====================================================================
% PART V -- QUICK REFERENCE & COMMON PATTERNS
% ====================================================================

\hh{NumPy recipes for common architectural tasks}

\begin{tcolorbox}[
    title=Copy-Paste Recipes,
    colframe=gray,
    colback = white,
    arc=2mm
]
\textbf{Recipe 1: Calculate FAR (Floor Area Ratio)}
\begin{lstlisting}[basicstyle=\ttfamily\footnotesize,breaklines=true,columns=fullflexible,keepspaces=true,xleftmargin=0pt]
# Given floor areas for each level
floor_areas = np.array([850, 750, 750, 750, 600, 600])
site_area = 2000  # m^2
total_gfa = floor_areas.sum()
far = total_gfa / site_area
print(f"FAR: {far:.2f}")

\end{lstlisting}

\textbf{Recipe 2: Generate Structural Grid}
\begin{lstlisting}[basicstyle=\ttfamily\footnotesize,breaklines=true,columns=fullflexible,keepspaces=true,xleftmargin=0pt]
# Create column grid points
x_spacing = 7.5  # meters
y_spacing = 8.0  # meters
x_bays = 8  # number of bays in X
y_bays = 5  # number of bays in Y

x_coords = np.arange(0, (x_bays + 1) * x_spacing, x_spacing)
y_coords = np.arange(0, (y_bays + 1) * y_spacing, y_spacing)
X, Y = np.meshgrid(x_coords, y_coords)
column_points = np.column_stack([X.flatten(), Y.flatten()])

\end{lstlisting}

\textbf{Recipe 3: Analyze Unit Mix}
\begin{lstlisting}[basicstyle=\ttfamily\footnotesize,breaklines=true,columns=fullflexible,keepspaces=true,xleftmargin=0pt]
# Count unit types and calculate percentages
unit_types = np.array(['studio', 'one-bed', 'one-bed', 
                       'two-bed', 'studio', 'three-bed'])
unique_types, counts = np.unique(unit_types, return_counts=True)
percentages = (counts / len(unit_types)) * 100

for unit_type, count, pct in zip(unique_types, counts, percentages):
    print(f"{unit_type}: {count} units ({pct:.1f}%)")

\end{lstlisting}

\textbf{Recipe 4: Calculate Facade Solar Exposure}
\begin{lstlisting}[basicstyle=\ttfamily\footnotesize,breaklines=true,columns=fullflexible,keepspaces=true,xleftmargin=0pt]
# For panels facing different directions
orientations = np.array([0, 45, 90, 135, 180, 225, 270, 315])  # degrees
sun_azimuth = 180  # South
sun_altitude = 45  # degrees

# Calculate incident angle (simplified)
angle_diff = np.abs(orientations - sun_azimuth)
angle_diff = np.minimum(angle_diff, 360 - angle_diff)
exposure_factor = np.cos(np.deg2rad(angle_diff))
exposure_factor = np.maximum(exposure_factor, 0)  # No negative

\end{lstlisting}

\textbf{Recipe 5: Net-to-Gross Ratio}
\begin{lstlisting}[basicstyle=\ttfamily\footnotesize,breaklines=true,columns=fullflexible,keepspaces=true,xleftmargin=0pt]
# Given room areas and circulation percentages
room_areas = np.array([45, 52, 48, 55, 42, 38])
circulation_factor = 1.25  # 25% for circulation
gross_areas = room_areas * circulation_factor
net_to_gross = room_areas.sum() / gross_areas.sum()
print(f"Net-to-gross ratio: {net_to_gross:.2%}")

\end{lstlisting}
\end{tcolorbox}

\hh{Quick Reference Card}

\begin{tcolorbox}[
    title=NumPy Essentials for Architects,
    colframe=gray,
    colback = white,
    arc=2mm
]
\textbf{Creating Arrays}
\begin{itemize}
\item From list: \c{np.array([1, 2, 3])}
\item Range: \c{np.arange(0, 10, 2)} gives [0, 2, 4, 6, 8]
\item Even spacing: \c{np.linspace(0, 10, 5)} gives 5 points from 0 to 10
\item Zeros: \c{np.zeros((3, 4))} gives 3x4 array of zeros
\item Ones: \c{np.ones((2, 5))} gives 2x5 array of ones
\end{itemize}

\textbf{Array Properties}
\begin{itemize}
\item \c{arr.shape} gives dimensions (rows, cols)
\item \c{arr.size} gives total number of elements
\item \c{arr.ndim} gives number of dimensions
\item \c{arr.dtype} gives data type of elements
\end{itemize}

\textbf{Selecting Data}
\begin{itemize}
\item Single element: \c{arr[0]} or \c{arr[2, 3]}
\item Row: \c{arr[2, :]} gives entire row 2
\item Column: \c{arr[:, 3]} gives entire column 3
\item Range: \c{arr[1:3, 0:2]} gives rows 1-2, columns 0-1
\end{itemize}

\textbf{Filtering}
\begin{itemize}
\item Create mask: \c{mask = arr > 50}
\item Apply mask: \c{arr[mask]}
\item Multiple conditions: \c{(arr > 30) \& (arr < 60)}
\item Find positions: \c{np.where(arr > 50)}
\end{itemize}

\textbf{Math Operations}
\begin{itemize}
\item All apply element-wise: \c{+}, \c{-}, \c{*}, \c{/}, \c{**}
\item Sum: \c{arr.sum()} or \c{arr.sum(axis=0)}
\item Mean: \c{arr.mean()}
\item Min/Max: \c{arr.min()}, \c{arr.max()}
\item Round: \c{np.round(arr, 2)}
\end{itemize}

\textbf{Reshaping}
\begin{itemize}
\item Reshape: \c{arr.reshape(3, 4)}
\item Flatten: \c{arr.flatten()}
\item Transpose: \c{arr.T}
\end{itemize}

\textbf{Axis Memory Aid}
\begin{itemize}
\item \c{axis=0} operates down (vertically, through the rows)
\item \c{axis=1} operates across (horizontally, through the columns)
\end{itemize}
\end{tcolorbox}

\hh{Common Patterns that we use often}

\begin{tcolorbox}[
    title=Patterns You'll Use Often,
    colframe=gray,
    colback = white,
    arc=2mm
]
\textbf{1. Calculate areas for all rooms:}
\begin{lstlisting}[basicstyle=\ttfamily\footnotesize,breaklines=true,columns=fullflexible,keepspaces=true,xleftmargin=0pt]
widths = np.array([...])
lengths = np.array([...])
areas = widths * lengths

\end{lstlisting}

\textbf{2. Find rooms meeting code requirements:}
\begin{lstlisting}[basicstyle=\ttfamily\footnotesize,breaklines=true,columns=fullflexible,keepspaces=true,xleftmargin=0pt]
min_area = 9.0  # m^2 minimum
compliant = areas >= min_area
compliant_rooms = areas[compliant]

\end{lstlisting}

\textbf{3. Calculate totals by floor:}
\begin{lstlisting}[basicstyle=\ttfamily\footnotesize,breaklines=true,columns=fullflexible,keepspaces=true,xleftmargin=0pt]
# building is 2D array (floors x rooms)
floor_totals = building.sum(axis=1)  # Sum across each floor

\end{lstlisting}

\textbf{4. Apply safety factors:}
\begin{lstlisting}[basicstyle=\ttfamily\footnotesize,breaklines=true,columns=fullflexible,keepspaces=true,xleftmargin=0pt]
loads = np.array([...])
factored = loads * 1.5

\end{lstlisting}

\textbf{5. Generate grid points:}
\begin{lstlisting}[basicstyle=\ttfamily\footnotesize,breaklines=true,columns=fullflexible,keepspaces=true,xleftmargin=0pt]
x = np.arange(0, width, spacing)
y = np.arange(0, height, spacing)
X, Y = np.meshgrid(x, y)
points = np.column_stack([X.flatten(), Y.flatten()])

\end{lstlisting}

\textbf{6. Snap to construction grid:}
\begin{lstlisting}[basicstyle=\ttfamily\footnotesize,breaklines=true,columns=fullflexible,keepspaces=true,xleftmargin=0pt]
raw_dims = np.array([3.142, 4.678, 5.823])
snapped = np.round(raw_dims, 1)  # 0.1m grid
snapped = np.round(raw_dims * 2) / 2  # 0.5m grid

\end{lstlisting}
\end{tcolorbox}

\hh{Debugging Tips}

\begin{tcolorbox}[
    title=When Things Go Wrong,
    colframe=gray,
    colback = white,
    arc=2mm
]
\textbf{Problem: "could not broadcast"}
\begin{itemize}
\item Check shapes: \c{print(arr1.shape, arr2.shape)}
\item Arrays must have compatible dimensions
\item One dimension must be 1 or they must be equal
\end{itemize}

\textbf{Problem: Unexpected results after slicing}
\begin{itemize}
\item You might be editing a view, not a copy
\item Use \c{.copy()} to ensure independence
\item Check with: \c{arr.base is not None} (True = it's a view)
\end{itemize}

\textbf{Problem: Wrong dimensions after reshape}
\begin{itemize}
\item Total elements must stay the same
\item 12 elements can be 3x4, 2x6, 12x1, but not 3x5
\item Use -1 to let NumPy figure out one dimension
\end{itemize}

\textbf{Problem: Can't find elements}
\begin{itemize}
\item \c{np.where()} returns indices, not values
\item Use: \c{indices = np.where(condition)[0]}
\item Then: \c{values = arr[indices]}
\end{itemize}
\end{tcolorbox}